\documentclass[11pt,a4paper]{article}
\usepackage[margin=2.5cm]{geometry}
\usepackage[T1]{fontenc}
\usepackage{booktabs}
\usepackage{xcolor}
\usepackage{amsmath}
\usepackage{listings}
\usepackage[hidelinks]{hyperref}

\emergencystretch=2em
\definecolor{codebg}{gray}{0.95}
\lstset{language=C++, basicstyle=\ttfamily\small, keywordstyle=\color{blue!70!black},
  commentstyle=\color{green!50!black}\itshape, backgroundcolor=\color{codebg},
  frame=single, rulecolor=\color{black!20}, breaklines=true, showstringspaces=false,
  tabsize=2, columns=flexible}

\title{ATmega32A Module \\ \Large PWM (Pulse Width Modulation)}
\author{Ali Sahafi}
\date{}

\begin{document}
\maketitle

\section{Theory}

PWM produces an ``analog-like'' output from a digital pin by switching it on
and off very fast. The fraction of time the pin is HIGH is the \textbf{duty
cycle}: at 25\,\% duty an LED appears dim, at 100\,\% fully bright. The
switching is far too fast for the eye (or a motor) to notice --- only the
\emph{average} voltage matters:

\begin{equation*}
V_{\text{avg}} = V_{\text{CC}} \cdot \frac{\text{duty}}{255}
\end{equation*}

In \textbf{Fast PWM} mode the timer counts 0--255 continuously; the output
pin is HIGH while the counter is below the compare register and LOW above
it. Writing the compare register therefore \emph{is} setting the duty cycle.
With the driver's fixed prescaler of 8:

\begin{equation*}
f_{\text{PWM}} = \frac{F_{\text{CPU}}}{8 \cdot 256} \approx 3.9\,\text{kHz at } 8\,\text{MHz}
\end{equation*}

PWM uses dedicated output pins --- the compare hardware drives them directly:

\begin{center}
\begin{tabular}{llll}
\toprule
\textbf{Channel} & \textbf{Pin} & \textbf{Timer} & \textbf{Duty register} \\
\midrule
OC0  & PB3 & Timer0 & \texttt{OCR0}  \\
OC1A & PD5 & Timer1 & \texttt{OCR1A} \\
OC1B & PD4 & Timer1 & \texttt{OCR1B} \\
\bottomrule
\end{tabular}
\end{center}

\section{API reference}

\begin{center}
\small
\begin{tabular}{ll}
\toprule
\textbf{Method} & \textbf{Description} \\
\midrule
\texttt{PWM.initTimer0\_FastPWM()} & Fast-PWM on PB3 (pin set to output automatically) \\
\texttt{PWM.setDutyCycleTimer0(duty)} & Set Timer0 duty cycle, 0--255 \\
\texttt{PWM.initTimer1\_FastPWM\_8bit()} & Fast-PWM on PD4 (OC1B) and PD5 (OC1A) \\
\texttt{PWM.setDutyCycleTimer1A(duty)} & Set PD5 duty cycle, 0--255 \\
\texttt{PWM.setDutyCycleTimer1B(duty)} & Set PD4 duty cycle, 0--255 \\
\bottomrule
\end{tabular}
\end{center}

\paragraph{Warning: Timer conflict.} PWM occupies the timer hardware, so the
CTC functions of Module 3 and the PWM functions must never share a timer:
\begin{itemize}
  \item \texttt{Timer.initTimer0()} conflicts with \texttt{PWM.initTimer0\_FastPWM()}
  \item \texttt{Timer.initTimer1()} conflicts with \texttt{PWM.initTimer1\_FastPWM\_8bit()}
\end{itemize}
Both write the same \texttt{TCCR} registers, and the later call silently
breaks the earlier one. Timer2 remains free for CTC interrupts alongside PWM.

\section{Example: breathing LED}

\paragraph{Wiring.}
\begin{itemize}
  \item LED + 330\,$\Omega$ resistor from \textbf{PB3} (OC0) to GND.
\end{itemize}

\paragraph{Build \& flash.} \texttt{make pwm}

\lstinputlisting{example.cpp}

Note the loop variables: counting \emph{up} uses \texttt{uint16\_t} and
counting \emph{down} uses a signed \texttt{int16\_t} --- an 8-bit
\texttt{uint8\_t} could never exit either loop. Why?

\section{Exercises}

\begin{enumerate}
  \item The brightness change looks ``front-loaded'': most of the visible
        change happens at low duty values. Why? (Hint: the human eye responds
        logarithmically.) Try stepping through a table of perceptually-spaced
        duty values instead.
  \item Use \texttt{initTimer1\_FastPWM\_8bit()} to fade two LEDs on PD4 and
        PD5 in \emph{opposite} directions (one brightens while the other dims).
  \item Connect the potentiometer from the ADC module (Module 5) and set the
        LED brightness from the potentiometer position: 10-bit ADC value
        $\rightarrow$ 8-bit duty cycle.
\end{enumerate}

\end{document}
